\section{La idea central de DDIM: un proceso forward no markoviano}

\subsection{Repaso del proceso forward de DDPM}

En DDPM, el proceso forward se define como una \textbf{cadena de Markov}:

$$
q(\mathbf{x}_{1:T} | \mathbf{x}_0) = \prod_{t=1}^{T} q(\mathbf{x}_t | \mathbf{x}_{t-1})
$$

donde la distribución de transición en cada paso es:
$$
q(\mathbf{x}_t | \mathbf{x}_{t-1}) = \mathcal{N}(\mathbf{x}_t; \sqrt{1-\beta_t}\,\mathbf{x}_{t-1},\, \beta_t \mathbf{I})
$$

De aquí se pueden deducir dos distribuciones clave:
\begin{itemize}
    \item \textbf{Distribución marginal}: $q(\mathbf{x}_t | \mathbf{x}_0) = \mathcal{N}(\mathbf{x}_t; \sqrt{\bar{\alpha}_t}\,\mathbf{x}_0, (1-\bar{\alpha}_t)\mathbf{I})$
    \item \textbf{Distribución posterior}: $q(\mathbf{x}_{t-1} | \mathbf{x}_t, \mathbf{x}_0)$, también una distribución gaussiana (solución en forma cerrada)
\end{itemize}

\begin{importantbox}{Las distribuciones realmente necesarias en el ELBO}
Al deducir la cota inferior variacional (ELBO), lo que necesitamos no es la propia distribución de transición forward $q(\mathbf{x}_t|\mathbf{x}_{t-1})$, sino:
\begin{enumerate}
    \item La distribución marginal $q(\mathbf{x}_t | \mathbf{x}_0)$, utilizada para muestrear los datos de entrenamiento.
    \item La distribución posterior $q(\mathbf{x}_{t-1} | \mathbf{x}_t, \mathbf{x}_0)$, objetivo de ajuste del proceso inverso.
\end{enumerate}
Esta observación es el punto de partida clave de DDIM: dado que el ELBO solo depende de estas dos distribuciones, ¿podemos saltarnos la definición de la distribución de transición forward y definir directamente estas dos?
\end{importantbox}

\subsection{La idea de DDIM: definir directamente la distribución marginal y la posterior}

La innovación central de DDIM radica en \textbf{invertir el orden de deducción}:

\begin{table}[H]
\centering
\caption{Comparación del orden de deducción entre DDPM y DDIM}
\begin{tabular}{ll}
\toprule
\textbf{DDPM} & \textbf{DDIM} \\
\midrule
Definir primero $q(\mathbf{x}_t|\mathbf{x}_{t-1})$ & Definir directamente $q(\mathbf{x}_t|\mathbf{x}_0)$ \\
Deducir $q(\mathbf{x}_t|\mathbf{x}_0)$ & Definir directamente $q(\mathbf{x}_{t-1}|\mathbf{x}_t, \mathbf{x}_0)$ \\
Deducir $q(\mathbf{x}_{t-1}|\mathbf{x}_t, \mathbf{x}_0)$ & La distribución de transición forward queda implícitamente determinada por los dos anteriores \\
\bottomrule
\end{tabular}
\end{table}

Concretamente, DDIM establece lo siguiente:

\paragraph{Mantener la distribución marginal invariable.} Para todo $t$, DDIM requiere:
$$
q(\mathbf{x}_t | \mathbf{x}_0) = \mathcal{N}(\mathbf{x}_t; \sqrt{\bar{\alpha}_t}\,\mathbf{x}_0,\, (1-\bar{\alpha}_t)\mathbf{I})
$$
Esto es exactamente lo mismo que en DDPM.

\paragraph{Redefinir la distribución posterior.} DDIM parametriza la distribución posterior como:
$$
q_\sigma(\mathbf{x}_{t-1} | \mathbf{x}_t, \mathbf{x}_0) = \mathcal{N}\left(\mathbf{x}_{t-1};\, \sqrt{\bar{\alpha}_{t-1}}\,\mathbf{x}_0 + \sqrt{1-\bar{\alpha}_{t-1}-\sigma_t^2}\cdot\frac{\mathbf{x}_t - \sqrt{\bar{\alpha}_t}\,\mathbf{x}_0}{\sqrt{1-\bar{\alpha}_t}},\, \sigma_t^2\mathbf{I}\right)
$$

Donde $\sigma_t$ es un hiperparámetro libre de elección que controla el grado de ruido inyectado en cada paso del proceso inverso.

\begin{warningbox}{El significado de la no markovianidad}
En el proceso forward de DDIM, la muestra de $\mathbf{x}_{t-1}$ depende de \textbf{ambos} $\mathbf{x}_t$ y $\mathbf{x}_0$ (no solo de $\mathbf{x}_t$). Esto rompe la hipótesis de Markov: el estado actual no solo depende del paso anterior, sino también de los datos originales. De aquí proviene el término ``Implicit'': la distribución de transición forward $q(\mathbf{x}_t|\mathbf{x}_{t-1})$ ya no se define explícitamente, sino que queda implícitamente determinada por la distribución marginal y la posterior.
\end{warningbox}

\subsection{Proceso de deducción: determinar coeficientes mediante inducción}

El ponente detalló cómo determinar los coeficientes de la distribución posterior a través del método de inducción matemática.

\paragraph{Restricción central.} Necesitamos cumplir que, para todo $t$, la $q(\mathbf{x}_{t-1}|\mathbf{x}_0)$ derivada mediante la distribución posterior $q(\mathbf{x}_{t-1}|\mathbf{x}_t, \mathbf{x}_0)$ y la distribución marginal $q(\mathbf{x}_t|\mathbf{x}_0)$ sea consistente con la distribución marginal de DDPM.

\paragraph{Composición en forma cerrada de distribuciones gaussianas.} El ponente aprovechó un hecho matemático clave:

\begin{knowledgebox}{Propiedad de marginalización de distribuciones gaussianas}
Si la distribución previa $p(\mathbf{z})$ y la distribución de verosimilitud $p(\mathbf{x}|\mathbf{z})$ son ambas gaussianas, y la media de la verosimilitud es una función lineal de $\mathbf{z}$, entonces la distribución marginal $p(\mathbf{x}) = \int p(\mathbf{x}|\mathbf{z})\,p(\mathbf{z})\,d\mathbf{z}$ también es gaussiana, y su media y varianza pueden calcularse mediante una expresión en forma cerrada. Este es el fundamento matemático que permite realizar toda la deducción.
\end{knowledgebox}

Asumiendo que la media de la distribución posterior es una combinación lineal de $\mathbf{x}_0$ y $\mathbf{x}_t$:
$$
\mu_{t-1}(\mathbf{x}_t, \mathbf{x}_0) = w_0 \cdot \mathbf{x}_0 + w_t \cdot \mathbf{x}_t
$$

Igualando la media y la varianza de $q(\mathbf{x}_{t-1}|\mathbf{x}_0)$, se obtiene:
$$
w_0 = \sqrt{\bar{\alpha}_{t-1}} - \sqrt{\bar{\alpha}_t} \cdot \frac{\sqrt{1-\bar{\alpha}_{t-1}-\sigma_t^2}}{\sqrt{1-\bar{\alpha}_t}}, \quad
w_t = \frac{\sqrt{1-\bar{\alpha}_{t-1}-\sigma_t^2}}{\sqrt{1-\bar{\alpha}_t}}
$$

\begin{warningbox}{Restricción de varianza: $\sigma_t$ no puede ser arbitrariamente grande}
Dado que la varianza debe ser no negativa, $\sigma_t^2$ tiene un límite superior:
$$
\sigma_t^2 \leq 1 - \bar{\alpha}_{t-1}
$$
Cuando $\sigma_t^2$ alcanza su límite superior, la parte de varianza de la distribución posterior coincide exactamente con la distribución posterior de DDPM. Esto indica que DDPM es un caso particular de DDIM bajo una elección específica de varianza.
\end{warningbox}

\subsection{Invariancia del ELBO}

\begin{importantbox}{El objetivo de entrenamiento permanece inalterado: no se requiere reentrenamiento}
Dado que DDIM mantiene la distribución marginal $q(\mathbf{x}_t|\mathbf{x}_0)$ idéntica a la de DDPM, sus ELBOs (cotas inferiores variacionales) son equivalentes salvo un factor constante independiente de los parámetros del modelo. Esto implica que:
\begin{enumerate}
    \item La red de predicción de ruido $\boldsymbol{\epsilon}_\theta$ entrenada con el objetivo de DDPM puede usarse directamente en el proceso de muestreo de DDIM.
    \item No se requiere ningún ajuste fino adicional ni reentrenamiento.
    \item Se logra aceleración únicamente mediante la modificación de la estrategia de muestreo durante la inferencia.
\end{enumerate}
Esta es la característica más valiosa en términos prácticos de DDIM: \textbf{entrenar una vez, razonar de múltiples maneras}.
\end{importantbox}

Este fenómeno es bastante raro en el aprendizaje profundo; usualmente, después de entrenar un modelo, la forma de razonación queda fija. Sin embargo, para los modelos difusos, debido a su estructura especial de muestreo iterativo, podemos elegir flexiblemente estrategias de muestreo durante la fase de inferencia (markoviana vs no markoviana, estocástica vs determinista) sin modificar el proceso de entrenamiento.

\subsection{Resumen de la sección}

La idea central de DDIM es: definir directamente la distribución marginal y la posterior (en lugar de la distribución de transición forward), construyendo así un proceso forward no markoviano. Mediante inducción matemática y las propiedades de composición en forma cerrada de distribuciones gaussianas, se pueden deducir expresiones analíticas para los coeficientes de la distribución posterior. Lo fundamental es que la distribución marginal de DDIM coincide completamente con la de DDPM; por lo tanto, el ELBO permanece invariable y un modelo de DDPM ya entrenado puede usarse directamente para el muestreo de DDIM.
